\documentclass[11pt,a4paper,twoside,openright,titlepage,final]{article}
\usepackage{fontspec}
\usepackage{amsmath}
\usepackage{amsfonts}
\usepackage{amssymb}
\usepackage{makeidx}
\usepackage{graphicx}
\usepackage[hidelinks,unicode=true]{hyperref}
\usepackage[spanish,es-nodecimaldot,es-lcroman,es-tabla,es-noshorthands]{babel}
\usepackage[left=3cm,right=2cm, bottom=4cm]{geometry}
\usepackage{natbib}
\usepackage{microtype}
\usepackage{ifdraft}
\usepackage{verbatim}
\usepackage[obeyDraft]{todonotes}
\ifdraft{
	\usepackage{draftwatermark}
	\SetWatermarkText{BORRADOR}
	\SetWatermarkScale{0.7}
	\SetWatermarkColor{red}
}{}
\usepackage{booktabs}
\usepackage{longtable}
\usepackage{calc}
\usepackage{array}
\usepackage{caption}
\usepackage{subfigure}
\usepackage{footnote}
\usepackage{url}
\usepackage{tikz}
\usepackage{pdflscape}
\usepackage{minted}

%\setsansfont[Ligatures=TeX]{texgyreadventor}
%\setmainfont[Ligatures=TeX]{texgyrepagella}

%*******************************************************
%                 NO MODIFICAR
\newcommand*{\FSfont}[1]{%
  \fontencoding{T1}\fontfamily{#1}\selectfont}

\newlength{\tpheight}\setlength{\tpheight}{0.9\textheight}
\newlength{\txtheight}\setlength{\txtheight}{0.9\tpheight}
\newlength{\tpwidth}\setlength{\tpwidth}{0.9\textwidth}
\newlength{\txtwidth}\setlength{\txtwidth}{0.9\tpwidth}
\newlength{\drop}
%*******************************************************

% Crea una portada con los siguientes parámetros
%
% #1 : Título 
% #2 : Subtítulo
% #3 : Subsubtítulo
% #4 : Autor(es)
% #5 : Lugar
%

\newcommand*{\portada}[5]{
\begin{titlepage}
\begingroup
\vspace*{1cm}
\drop = 0.2\txtheight
\centering
\vfill
{\Huge \scshape #1}\\[\baselineskip]
{\Large \textbf{#2}}\\[\baselineskip]
{\Large \scshape #3}\\[\baselineskip]
\vspace*{0.3cm}
{\large \textit{#4}}\\[0.5\drop]
\includegraphics[scale=0.35]{./imagenes/logoURJC.jpg}
\vspace*{1.5cm}

{\large \scshape #5, \today} \par
\begin{center}
\end{center}
\vfill\null
\endgroup
\end{titlepage}
}
 %*****************************************************
 


\author{José Ignacio Escribano}

\setlength{\parindent}{0pt}

\begin{document}

\pagenumbering{alph}
\setcounter{page}{1}

\portada{Práctica I}{Simulación y Metaheurísticas}{Problemas e instancias}{José Ignacio Escribano}{Móstoles}

\listoftables
\thispagestyle{empty}
\newpage

\tableofcontents
\thispagestyle{empty}
\newpage


\pagenumbering{arabic}
\setcounter{page}{1}

\section{Introducción}

El problema del capacited p-hub consiste en determinar los $p$ centros (hubs) que actúan como servidores y conectar todos los clientes a uno de los servidores de forma que se minimice el la suma de las distancias de los clientes a los servidores. Además, cada servidor sólo puede puede servir una determinada cantidad de recursos, y cada cliente tiene una demanda que deben ser satisfechas. De forma matemática, el modelo se formula así:

\begin{alignat*}{2}
\text{Minimizar }   & \sum_{i} \sum_{j} d_{ij}x_{ij}  \\
\text{s.a. } &  \sum_{j} x_{ij} = 1 \quad \forall i\\
&  x_{ij} \leq y_j \quad \forall i,j\\
& \sum_{j} y_j = p \\
& x_{ij}, y_j \in \{0,1\} \\
& \sum_{i} x_{ij} b_j \leq c
\end{alignat*}

donde

\begin{itemize}
	\item $x_{ij}$ es $1$ si hay arista entre el cliente $i$ y el servidor $j$.
	\item $y_i$ es $1$  si $j$ es servidor, $0$ en caso contrario.
	\item $d_{ij}$ es la distancia entre el nodo $i$ y el $j$.
	\item $b_j$ es la demanda del nodo $j$
	\item $c$ es la capacidad del servidor
\end{itemize}

\section{Resolución de la práctica}

A continuación, resolveremos las cuestiones planteadas en la práctica.

\subsection{Estructuras de datos de las instancias}

Las instancias para el problema del \textit{p-hub} se encuentran en la clase \texttt{InstanciaPHub}, cuyos atributos son los siguientes:

\begin{enumerate}
	\item nodos: número de nodos de la instancia.
	\item servidores: número de servidores de la instancia (p).
	\item distancia: matriz cuadrada que contiene la distancia entre cada par de nodos, calculada a partir de las coordenadas de cada centro usando la distancia euclídea, esto es,
	
	\begin{equation}
	\label{eq:distancia}
	d(\mathbf{x} =(x_1,\dots, x_n), \mathbf{y} = (y_1,\dots, y_n)) = \sqrt{\sum_{i=1}^{n} (x_i - y_i)^2 }
	\end{equation}
	
	\item demanda: vector con la demanda de cada nodo.
	\item capacidad: capacidad de los servidores.
\end{enumerate}

\subsection{Carga de las instancias}

Las instancias se cargan en memoria gracias al procedimiento \texttt{leerInstancia} de la clase \texttt{PHub}. Este procedimiento recibe como parámetro el nombre del fichero y va leyendo línea a línea los datos del fichero. Estos datos se guardan en sus respectivas estructuras de datos, que se pasan al constructor de la clase \texttt{InstanciaPHub}, que ``genera'' el objeto. Este procedimiento también se encarga de almacenar las posiciones de los centros y las almacena en memoria, para llamar a la función \texttt{calcularDistancias} de la clase \texttt{PHub}, que se encarga de devolver una matriz de distancias haciendo uso de la Ecuación~\ref{eq:distancia}.

\subsection{Estructura de la solución}

Cada solución se guarda en un objeto de la clase \texttt{Solución}, que contiene los siguientes atributos:

\begin{enumerate}
	\item solución: vector de booleanos que indica \texttt{true}, si el nodo $i$ es un servidor, \texttt{false} en caso contrario.
	\item matrizAdyacencia: matriz de booleanos, que contiene la matriz de adyacencia del grafo, esto es, \texttt{true}, si hay una arista entre el nodo $i$ y $j$, \texttt{false} en caso contrario.
	\item objetivo: valor de la función objetivo, es decir, la suma de las aristas existentes entre los clientes y los objetivos. 
\end{enumerate}

\subsection{Cálculo de la función objetivo}

El cálculo de la función objetivo se realiza en la función \texttt{calcularObjetivo} de la clase \texttt{Utils}, que dada una solución, su matriz de adyacencia y de distancias, calcula su valor.

El código es el siguiente:

\begin{minted}[linenos=true, numberblanklines=true, breaklines=true]{java}
public static double calcularObjetivo(boolean[] solucion, boolean[][] matrizAdyacencia, double[][] distancia) {

    double obj = 0;
    for (int i = 0; i < matrizAdyacencia.length; i++) {
        for (int j = 0; j <= i; j++) {
            // Si hay conexión entre los nodos, sumamos
            if (matrizAdyacencia[i][j]) {
                obj += distancia[i][j];
            }
        }
    }
    return obj;
}
\end{minted}

Esta función recorre la matriz de adyacencia (notar que es simétrica), y si existe una arista entre el nodo $i$ y el $j$, suma la distancia entre estos dos nodos a la variable \texttt{obj}, que se devuelve.

\subsection{Solución aleatoria de una instancia}

Para generar una solución aleatoria de una instancia está la función \texttt{generarSoluciónAleatoria} de la clase \texttt{InstanciaPHub}. El código es el siguiente:

\begin{minted}[linenos=true, numberblanklines=true, breaklines=true]{java}
public Solución generarSoluciónAleatoria() {

    Solución s1 = null;

    Random x = new Random();

    // Para poder reproducir los resultados
    x.setSeed(0);

    do{
        x = new Random();

        int nodos = this.getNodos();
        int num_servidores = 0;

        boolean[] sol = new boolean[nodos];
        boolean[][] ady = new boolean[nodos][nodos];

        // Elegimos al azar los servidores
        while (num_servidores < this.getServidores()) {
            int aleatorio = 0;
        
            do {
                aleatorio = x.nextInt(nodos);
            } while (sol[aleatorio]);

            sol[aleatorio] = true;

            num_servidores++;
        }

        // Generamos las aristas
        for (int i = 0; i < nodos; i++) {
            if (!sol[i]) {
                // Seleccionamos el servidor más cercano que nos encontremos
                int serv = Utils.seleccionarServidor(i, sol, this.getDistancia());

                ady[i][serv] = true;
                ady[serv][i] = true;
            }

         }

         s1 = new Solución(sol, ady, this.getDistancia());

    }while(!Utils.esSoluciónVálida(s1, this));

    return s1 ;
}
\end{minted}

En las líneas 20 a 30 se seleccionan al azar los $p$ nodos que actuarán como servidores. Para ello, se generan números aleatorios en el intervalo discreto $[0, \text{número de nodos})$. Una vez que se tienen los nodos que actuarán como servidores, se selecciona a qué servidor se conectará cada cliente. Para ello, se busca el servidor más cercano (el que tiene una menor distancia euclídea) al cliente. De esto se encarga la función \texttt{seleccionarServidor} de la clase \texttt{Utils}. Una vez se tiene el servidor más cercano al cliente, se establece a \texttt{true} la posición de la matriz de adyacencia entre el servidor asignado y el cliente (notar que la matriz de adyacencia es simétrica). Con todos estos datos, se crea un objeto de la clase \texttt{Solución} y se comprueba si es una solución válida, es decir, si cumple las restricciones del problema. Si lo es, se devuelve la solución, y en caso contrario se repite el proceso anterior hasta que se obtenga una solución válida.\\

El código de la función \texttt{esSoluciónVálida} de la clase \texttt{Utils} es el siguiente:

\begin{minted}[linenos=true, numberblanklines=true, breaklines=true]{java}
public static boolean esSoluciónVálida(Solución s, InstanciaPHub instancia) {
    boolean esValida = false;

    boolean[][] ady = s.getMatrizAdyacencia();
    boolean[] sol = s.getSolucion();
    int capacidad = instancia.getCapacidad();

    // Comprobamos si se cumple el criterio de la demanda
    for (int i = 0; i < sol.length; i++) {
        int demanda = 0;
        for (int j = 0; j < sol.length; j++) {
            if (sol[i] && ady[i][j]) {
                demanda += instancia.getDemanda()[j];
            }
        }

        if (demanda <= capacidad && demanda != 0) {
            esValida = true;
        } else {
	        esValida = false;
        }
     }
     return esValida;
}
\end{minted}

En las líneas 9 a 20 se recorre el vector solución (que contiene \texttt{true}, si el nodo $i$ es un servidor, y \texttt{false} en caso contrario) para comprobar si se cumple la restricción de la demanda, esto es, si la suma de las demanda de cada cliente es menor o igual a la capacidad del servidor. Si se comprueba que todos los servidores tienen una demanda menor a la capacidad se devuelve \texttt{true}. En caso contrario \texttt{false}.

\subsection{Programa principal}

El programa principal se encarga de generar $1000$ soluciones aleatorias de cada una de las instancias, y devolviendo el valor de la función objetivo y la estructura de la mejor solución de cada instancia.\\

Los valores de la mejor solución de cada instancia se muestran en la Tabla~\ref{tbl:mejores}.\\

\begin{table}[htbp!]
	\centering
	\caption{Mejores soluciones de cada instancia}
	\label{tbl:mejores}
	\begin{tabular}{@{}ccc@{}}
		\toprule
		Instancia & Función objetivo & Servidores               \\ \midrule
		1         & 768.543          & {[}5,9,10,21,23{]}       \\
		2         & 833.446          & {[}4,5,6,12,16{]}        \\
		3         & 802.602          & {[}22,24,36,38,48{]}     \\
		4         & 797.988          & {[}15,34,38,47,50{]}     \\
		5         & 709.911          & {[}3,7,8,11,23{]}        \\
		6         & 687.306          & {[}14,16,22,36,40{]}     \\
		7         & 851.012          & {[}11,19, 32,41,50{]}    \\
		8         & 800.853          & {[}4, 13, 32, 36, 48{]}  \\
		9         & 802.646          & {[}14, 18, 28, 40, 47{]} \\
		10        & 756.424          & {[}10, 12, 22, 23, 38{]} \\ \bottomrule
	\end{tabular}
\end{table}

Notar que se han omitido la matriz de adyacencia por claridad. Las soluciones completas se pueden encontrar en el fichero \texttt{mejores\_soluciones.txt} adjunto a esta memoria.


\section{Conclusiones}

En esta práctico nos hemos introducido en el problema p-hub. Hemos cargado instancias de este problema en memoria y hemos generado soluciones aleatorias de cada una de las instancias. En último lugar hemos simulado 1000 soluciones de cada instancia y hemos seleccionado la mejor (la que tenía menor valor de la función objetivo).

\end{document}